\vspace{-.15cm}
\section{Conclusion}
\label{sec:discussion}
We have presented \method, a denoising diffusion model for trajectory data.
Planning with \method is almost identical to sampling from it, differing only in the addition of auxiliary perturbation functions that serve to guide samples.
The learned diffusion-based planning procedure has a number of useful properties, including graceful handling of sparse rewards, the ability to plan for new rewards without retraining, and a temporal compositionality that allows it to produce out-of-distribution trajectories by stitching together in-distribution subsequences.
% Empirically, \method is particularly strong in control settings requiring long-horizon reasoning and test-time flexibility.
Our results point to a new class of diffusion-based planning procedures for deep model-based reinforcement learning.
% a new class of generative models and associated planning procedures for deep model-based reinforcement learning.
% encoding rewards
% that can be used for behavior synthesis.
% Unlike conventional single-step dynamics models, \method lends itself to a learned, non-greedy planning procedure that allows it so solve long-horizon tasks.
% Planning with \method amounts to sampling with auxiliary perturbation functions, allowing for test-time flexibility by encoding new reward functions in a new perturbation function and keeping the diffusion model unchanged.
% In this paper, we present \method{}, a probabilistic planning framework for flexible behavior synthesis.
% Our probabilistic planning framework exhibits unique properties of non-greedy planning, temporal compositionality, variable-length prediction and cost compositionality. 
% We illustrate how this enables \method{} to exhibit long horizon multi-task planning, test-time behavior flexibility, and more  effective offline reinforcement learning.
% Our results point to \method{} as a  promising direction of future direction of exploration for the generation of flexible and general purpose behavior.
% \todo{TODO}